\documentclass[11pt,a4paper]{article}

% ---- Page geometry per ICCE 2026 template (A4: top 1in, bottom 0.8in, left/right 1in) ----
\usepackage[a4paper,top=1in,bottom=0.8in,left=1in,right=1in]{geometry}

% ---- Arial-equivalent sans-serif body font, per template spec ----
\usepackage[T1]{fontenc}
\usepackage{helvet}
\renewcommand{\familydefault}{\sfdefault}

\usepackage[english]{babel}
\usepackage{ragged2e}
\usepackage{booktabs}
\usepackage{array}
\usepackage{caption}
\usepackage{graphicx}
\usepackage{amsmath}
\usepackage{changepage}
\usepackage{titlesec}
\usepackage{fancyhdr}
\usepackage{natbib}
\bibliographystyle{apalike}

% No headers/footers/page numbers per template instructions
\pagestyle{empty}

% ---- ICCE proceedings header, first page only ----
\fancypagestyle{iccefirst}{%
  \fancyhf{}%
  \renewcommand{\headrulewidth}{0pt}%
  \fancyhead[L]{\fontsize{9}{11}\selectfont Song, Y. et al. (Eds.) (2026). \textit{Proceedings of the 34th International Conference on Computers in Education}. Asia-Pacific Society for Computers in Education}%
}


% Single spacing, justified body text
\linespread{1.0}

% ---- Heading styles per template: first-level = bold, numbered, flush left;
% ---- sub-headings (2.1, 2.1.1, ...) = italic (not bold), flush left ----
\titleformat{\section}{\normalfont\fontsize{12}{14}\selectfont\bfseries}{\thesection.}{0.5em}{}
\titlespacing{\section}{0pt}{1.4\baselineskip}{0.7\baselineskip}
\titleformat{\subsection}{\normalfont\fontsize{11}{13}\selectfont\itshape}{\thesubsection}{0.5em}{}
\titlespacing{\subsection}{0pt}{0.7\baselineskip}{0.5\baselineskip}

% Table captions: roman label, italic caption text, above table, left-aligned, no full stop
\captionsetup[table]{position=top, justification=raggedright, singlelinecheck=false, labelfont=normalfont, textfont=it, labelsep=period}
% Figure captions: italic label, roman caption text, below figure, centered
\captionsetup[figure]{position=bottom, justification=centering, labelfont=it, textfont=normalfont, labelsep=period}

\setlength{\parindent}{1.25cm}
\setlength{\parskip}{0pt}

\begin{document}
\thispagestyle{iccefirst}
\justifying

% ---------------- TITLE BLOCK ----------------
\begin{center}
{\fontsize{22}{26}\selectfont\bfseries Temporal Sampling Strategies for Student Engagement Detection}

\vspace{0.35cm}

{\large Author(s) Anonymous}

\vspace{0.2cm}

{\itshape $^{a}$Anonymous Institution}

\vspace{0.15cm}

\end{center}

\vspace{0.3cm}

% ---------------- ABSTRACT ----------------
% Abstract style per template: 10pt, indented 0.5in (1.27cm) each side from the body margins
\begin{adjustwidth}{1.27cm}{1.27cm}
\fontsize{10}{12}\selectfont
\noindent\textbf{Abstract:} Automatic engagement recognition from video requires selecting a subset of frames before feature extraction, yet frame selection is rarely treated as a design variable. Four temporal sampling strategies were crossed with four classification architectures on EngageNet under a subject-independent split. The strategies were a Start-Middle-End (SME) positional baseline and three content-based methods selecting the top 3, 5, or 7 facial-landmark velocity peaks. The architectures were a Long Short-Term Memory network (LSTM), a Multi-Layer Perceptron (MLP), a Support Vector Machine (SVM), and a Random Forest. Every architecture reached its best macro F1 score under 3-Changepoint sampling, and performance declined as further frames were added. The best configuration, 3-Changepoint with SVM, reached 0.7010 macro F1. Paired McNemar testing shows the effect is not uniform. It reaches significance only for Random Forest, the weakest architecture, indicating that frame selection matters most where a classifier is least able to compensate for uninformative frames. Content-based sampling also reduced LSTM seed variance six-fold without improving its accuracy. Temporal sampling therefore warrants deliberate tuning, but its benefit depends on model capacity rather than applying uniformly across architectures.

\vspace{0.2cm}
\noindent\textbf{Keywords:} Student engagement, engagement recognition, temporal sampling, changepoint detection, facial action units
\end{adjustwidth}

\vspace{0.25cm}

% ---------------- 1. INTRODUCTION ----------------
\section{Introduction}

Student engagement is a well-established predictor of learning outcomes \citep{fredricks2004}. As online and hybrid learning environments have become more common, the loss of an instructor's ability to directly observe a room of students has made automatic, camera-based engagement detection more useful. Engagement signals detected from video also serve as behavioral inputs to agentic AI systems in educational contexts, where they can inform interventions while preserving learner agency.

Progress has depended on larger and more diverse datasets, and two now dominate the area. DAiSEE comprises 9,068 ten-second snippets from 112 participants, mostly university students aged 18 to 30, who filmed their own upper body with a webcam while watching massive open online course (MOOC) lectures in unconstrained settings such as dormitories and libraries \citep{gupta2016}. EngageNet was introduced to address the scale and annotation limits of corpora such as DAiSEE, following the same in-the-wild approach at a larger scale \citep{singh2023}. Its 127 participants, aged 18 to 37, watched three five-minute educational videos on quantum mechanics, cryptocurrency, and English language, each narrated by a digital avatar acting as a virtual teacher, and answered two comprehension questions after each video so that reported engagement could be checked against learning. Recordings were captured through a web platform under varying illumination, divided into roughly 11,300 ten-second clips, and labelled by three expert annotators.

Because engagement datasets consist of video rather than static images, any classification pipeline must first decide \textit{how} to reduce a clip to a tractable set of frames before features can be extracted. This step is often treated as a minor preprocessing detail. Prior engagement work has relied on fixed, position-based schemes, either sparse positional sampling such as taking the first, middle, and last frame (Start-Middle-End, or SME), or uniform segmentation, which divides a clip into equal parts and aggregates across them. Uniform segmentation is used in the DAiSEE engagement benchmark \citep{gupta2016} and in the EngageNet baselines \citep{singh2023}, and derives from the temporal segment network approach to action recognition \citep{wang2016}. Content-aware methods that instead select frames by motion salience have improved downstream accuracy in tasks such as action recognition \citep{korbar2019}. One way to identify informative frames by content is changepoint detection, a well-established family of methods for locating points in a signal where its underlying statistics shift \citep{truong2020}. Applied to facial-landmark velocity, this offers a content-based alternative to positional sampling that remains largely untested for engagement recognition specifically.

Does moving from position-based to content-based temporal sampling improve engagement classification performance across different model architectures? To address this, a factorial ablation crossed four temporal sampling strategies with four classification architectures on EngageNet under subject-independent binary classification. Two things were measured: whether content-based selection outperforms SME sampling consistently across architectures, and how the number of retained frames affects performance.

% ---------------- 2. METHODOLOGY ----------------
\section{Methodology}

\subsection{Dataset and Label Preprocessing}

A factorial ablation design was used, crossing sampling strategy (SME, 3-, 5-, and 7-Changepoint) with architecture (LSTM, MLP, SVM, Random Forest) to isolate the effect of each factor from the other. The data used was the EngageNet dataset, a collection of short video clips of students in online-learning activities, annotated per clip on a four-point engagement scale: 0 (Not-Engaged), 1 (Barely-Engaged), 2 (Engaged), and 3 (Highly-Engaged) \citep{singh2023}. Rows labelled SNP (Subject-Not-Present) carry no valid engagement score and were dropped prior to training.

The official EngageNet partition was used without re-partitioning. It is subject-independent, meaning no participant appears in more than one split, which excludes the failure mode in which a model identifies individual students rather than engagement states. Restricting every condition to the clips covered by all four sampling strategies (Section~2.3) reduces the 32-subject training split from 2,825 to 2,719 clips, the 11-subject validation split from 1,068 to 867, and the 26-subject test split from 2,227 to 2,005. At the binarization threshold used (Section~2.4), the restricted training split holds 1,375 negative and 1,344 positive samples, close to balanced.

\subsection{Feature Engineering Pipeline}

Per-frame features were extracted with MediaPipe, producing a 1,518-dimensional vector per frame from three families. The first family is 52 blendshape activation scores, which are continuous analogues of Facial Action Units (AU) \citep{ekman1978} rather than coder-assigned Action Unit intensities, capturing localised expressive movement such as brow raising, cheek squinting, and mouth-corner displacement. The second is 1,434 facial-landmark (FLM) values, comprising 478 landmarks in three dimensions, retaining geometric configuration not reduced to expression codes. The third is 32 pose keypoint (MP) values, comprising eight upper-body keypoints each with three coordinates and a visibility score, capturing postural and gestural information beyond the face.

Extraction proceeded in two passes. A first pass scanned every second frame to detect face and pose landmarks, establishing which frames were valid and producing the signal used for selection. A frame counted as valid when a face was detected and shoulder visibility was mutually consistent. A second pass then re-ran inference on the selected frames only.

\subsection{Temporal Sampling Strategies}

Four temporal sampling strategies were compared. The Start-Middle-End (SME) baseline retains the first, middle, and last validly-detected frame of a clip, choosing three frames by fixed position. The three Changepoint strategies instead retain the top-$k$ peaks in facial-landmark velocity ($k = 3, 5,$ or $7$), choosing frames by content rather than position, following the changepoint detection framework described by \citet{truong2020}. Both approaches sample sparsely and differ only in whether selection is driven by position or by content.

The changepoint procedure operated as follows. Facial-landmark coordinates from valid frames were mean-centred within each frame, and the Euclidean displacement between consecutive frames was taken as a velocity signal. That signal was smoothed with a five-frame moving average. Candidate peaks were located with a standard peak-finding routine constrained to a minimum separation of seven frames, which prevents several adjacent frames of one continuous movement from being selected as distinct peaks. No height or prominence threshold was applied, so the $k$ candidates with the largest velocity were retained regardless of their absolute magnitude, then restored to temporal order before feature extraction. When a clip contained fewer valid frames than $k$, all valid frames were kept. This last case is why the number of usable clips differs across strategies, and it motivates the common-subset restriction described in Section~2.1.

\subsection{Model Architectures and Binary Classification}

Four architectures were evaluated, spanning sequence-aware and pooled representations. The Long Short-Term Memory network (LSTM) received the ordered frame sequence through a 64-unit recurrent layer, then a 32-unit fully-connected layer, with a mean-pooled readout across timesteps. The Multi-Layer Perceptron (MLP) matched those widths but received average-pooled features, serving as the non-sequential control. The Support Vector Machine (SVM) used a radial basis function kernel on standard-scaled pooled features, and the Random Forest used 100 unpruned trees on the same input. All models used library defaults with no hyperparameter search, so comparisons are between architectures at default rather than best-tuned settings.

A binarization threshold of 3 was used, giving class 1 (C1), containing only Highly-Engaged clips, against class 0 (C0), which pools Not-Engaged, Barely-engaged, and Engaged. This isolates the class of practical interest and provided the most balanced distribution of the thresholds considered.

\subsection{Evaluation Metrics}

The following metrics were reported for each configuration: accuracy, macro F1, per-class F1, Pearson correlation coefficient with ground truth (PCC), mean squared error (MSE), and model training time. Macro F1 is unweighted across classes and is therefore sensitive to minority-class performance regardless of class size.

Because every condition is trained and evaluated on identical clips, predictions can be paired per clip, and exact McNemar tests were used to compare conditions. Two families were tested: each sampling strategy against the SME baseline within an architecture, and each architecture pair within a sampling condition. Holm-Bonferroni correction was applied within each family. The neural models were seeded and are reported as the mean over five seeds, whereas the SVM and Random Forest are deterministic given fixed data and so have no seed variance to report.

% ---------------- 3. RESULTS ----------------
\section{Results}

\subsection{Main Results}

Table~\ref{tab:main} reports accuracy, macro F1, per-class F1, PCC, MSE, and training time for all 16 sampling-strategy $\times$ architecture combinations at the binarization threshold described in Section~2.4, without oversampling ($n = 2{,}005$).

\begin{table}[h]
\caption{Accuracy and Related Metrics by Sampling Strategy and Architecture (Threshold 3, n = 2,005). C0 Is the Pooled Not-Highly-Engaged Class (Levels 0--2) and C1 Is Highly-Engaged (Level 3). LSTM and MLP Are Means $\pm$ SD Over Five Seeds. SVM and Random Forest Are Deterministic and Have No Seed Variance}
\label{tab:main}
\centering
\fontsize{9}{10}\selectfont
\begin{tabular}{llrrrrrrr}
\toprule
Sampling & Model & Accuracy & Macro F1 & C0 F1 & C1 F1 & PCC & MSE & Train (s) \\
\midrule
SME & LSTM & 0.6480$\pm$0.0129 & 0.6475 & 0.6526 & 0.6424 & 0.2983 & 0.3520 & 10.4 \\
SME & MLP & 0.6509$\pm$0.0096 & 0.6497 & 0.6485 & 0.6508 & 0.3042 & 0.3491 & 9.3 \\
SME & SVM & 0.6848 & 0.6838 & 0.7013 & 0.6663 & 0.3754 & 0.3152 & 2.6 \\
SME & RF & 0.5651 & 0.5307 & 0.6578 & 0.4036 & 0.1705 & 0.4349 & 1.8 \\
3-CP & LSTM & 0.6553$\pm$0.0022 & 0.6548 & 0.6511 & 0.6585 & 0.3113 & 0.3447 & 10.6 \\
3-CP & MLP & 0.6630$\pm$0.0047 & 0.6617 & 0.6450 & 0.6784 & 0.3267 & 0.3370 & 9.2 \\
\textbf{3-CP} & \textbf{SVM} & \textbf{0.7017} & \textbf{0.7010} & \textbf{0.7163} & \textbf{0.6856} & \textbf{0.4090} & \textbf{0.2983} & 2.7 \\
3-CP & RF & 0.6130 & 0.5904 & 0.6866 & 0.4941 & 0.2700 & 0.3870 & 3.0 \\
5-CP & LSTM & 0.6454$\pm$0.0077 & 0.6413 & 0.6143 & 0.6683 & 0.2945 & 0.3546 & 13.0 \\
5-CP & MLP & 0.6491$\pm$0.0045 & 0.6471 & 0.6333 & 0.6609 & 0.3001 & 0.3509 & 9.7 \\
5-CP & SVM & 0.6973 & 0.6966 & 0.7111 & 0.6820 & 0.3995 & 0.3027 & 2.5 \\
5-CP & RF & 0.5910 & 0.5612 & 0.6756 & 0.4467 & 0.2290 & 0.4090 & 2.0 \\
7-CP & LSTM & 0.6520$\pm$0.0080 & 0.6494 & 0.6336 & 0.6652 & 0.3068 & 0.3480 & 11.2 \\
7-CP & MLP & 0.6466$\pm$0.0121 & 0.6445 & 0.6299 & 0.6591 & 0.2950 & 0.3534 & 9.8 \\
7-CP & SVM & 0.6898 & 0.6884 & 0.7091 & 0.6677 & 0.3871 & 0.3102 & 2.6 \\
7-CP & RF & 0.5890 & 0.5610 & 0.6720 & 0.4499 & 0.2213 & 0.4110 & 2.1 \\
\bottomrule
\end{tabular}
\end{table}

The best-performing configuration overall was 3-Changepoint sampling with an SVM classifier (0.7010 macro F1, 0.6856 F1 on the Highly-Engaged class). Every architecture reached its own best macro F1 score under 3-Changepoint sampling rather than under the SME baseline, so the direction of the effect is consistent across architectures rather than specific to one model.

A consistent direction does not establish that the differences are large enough to be reliable. Table~\ref{tab:signif} reports paired McNemar tests of each changepoint condition against the SME baseline within each architecture.

\begin{table}[h]
\caption{Accuracy Difference From the SME Baseline in Percentage Points, With Holm-Bonferroni Corrected McNemar $p$-values. Positive Values Favour Changepoint Sampling}
\label{tab:signif}
\centering
\begin{tabular}{lrrr}
\toprule
Model & 3-CP & 5-CP & 7-CP \\
\midrule
Random Forest & $+4.79$ ($p<0.001$)* & $+2.59$ ($p=0.004$)* & $+2.39$ ($p=0.004$)* \\
SVM & $+1.70$ ($p=0.082$) & $+1.25$ ($p=0.209$) & $+0.50$ ($p=0.558$) \\
MLP & $+1.05$ ($p=0.413$) & $-0.45$ ($p=0.611$) & $-1.65$ ($p=0.233$) \\
LSTM & $-0.15$ ($p=0.890$) & $-2.49$ ($p=0.011$)*\textsuperscript{a} & $-0.60$ ($p=0.751$) \\
\bottomrule
\end{tabular}

\vspace{0.15cm}
{\fontsize{9}{10}\selectfont * Significant after correction. \textsuperscript{a}Seed artefact: McNemar uses one prediction set, and this difference reverses sign in three of the five seeds.}
\end{table}

Only the Random Forest shows a sampling effect that survives correction. The SVM gains 1.70 points under 3-Changepoint sampling, but this does not reach significance, and neither neural model shows a reliable effect. The benefit of content-based sampling therefore appears to depend on model capacity. The Random Forest, which is the weakest architecture by 8.9 to 12.0 accuracy points, gains the most, while the SVM gains least. A plausible reading is that frame selection matters more when the classifier cannot compensate for uninformative frames on its own, though the present design cannot separate this from other differences between the two model families.

\subsection{Effect of Frame Count}

Within the changepoint family, macro F1 declines as the frame budget grows. For the two deterministic models the decline from 3- to 7-Changepoint sampling is monotonic and noise-free, falling 1.26 points for the SVM and 2.94 points for the Random Forest. For the neural models the ordering is less clean and the differences sit within seed variation. This pattern is consistent with the extraction method itself. Because frames are selected by ranked velocity peak, frames beyond the top few are by construction weaker peaks, and so are more likely to capture blinks or camera jitter than engagement-relevant transitions. Controlled frame-sampling benchmarks report the same pattern, with larger frame budgets failing to improve performance monotonically and the best strategy varying by dataset and task \citep{brkic2025}. The present results indicate that this dependency also holds within a single dataset and task, at the level of how many content-selected frames are retained.

Sampling strategy also affects training stability, which is measurable only because the runs were seeded. Across five seeds, LSTM accuracy has a standard deviation of 0.0129 under SME sampling but only 0.0022 under 3-Changepoint sampling, a six-fold reduction, with the 5- and 7-Changepoint conditions intermediate at 0.0077 and 0.0080. The MLP shows the same effect more weakly, a roughly two-fold reduction. Content-based frame selection therefore improves reproducibility even where it does not improve accuracy.

Additional frames also increase extraction cost without improving accuracy. Extraction is dominated by a constant first-pass scan of roughly 4.5~s per 10~s clip, with re-inference adding about 0.9~s per selected frame. SME and 3-Changepoint sampling averaged 6.0--6.2~s per clip against 8.4~s for 7-Changepoint, a 7.8\% increase for no gain.

\subsection{Feature Importance}

Because the Random Forest discards frame ordering entirely yet is compared against sequence-aware models, its Gini feature importances were examined. Brow movement is the strongest single signal under both SME and 3-Changepoint sampling, with outer-brow raise ranking first and inner-brow raise in the top three of both. Because this holds under both strategies, it is more likely a property of the task than of one sampling configuration. Importance is also diffuse, the top feature carrying roughly 22 times its uniform share, so decisions rest on many weak signals rather than a few strong ones. Left-wrist visibility ranked unexpectedly high under both strategies, plausibly proxying hand-raising or note-taking, but separating genuine behavioural signal from a pose-detector artefact would require dedicated analysis.

% ---------------- 4. CONCLUSION ----------------
\section{Conclusion}

The hypothesis that content-based temporal sampling outperforms position-based approaches across different architectures was tested. Across all four models, 3-Changepoint sampling produced the best macro F1 score, and the best configuration, 3-Changepoint with an SVM, achieved 0.7010 macro F1. Paired McNemar testing qualifies that result. Only the Random Forest, the weakest of the four architectures, shows a sampling effect that survives correction, while the SVM gains least. The benefit of content-based selection therefore depends on model capacity rather than applying uniformly across architectures.

Adding more changepoints did not improve performance. Macro F1 declined monotonically from 3- to 7-Changepoint sampling for the two deterministic models while extraction grew more expensive, which is consistent with frames beyond the top few peaks capturing incidental motion such as blinks rather than engagement-relevant transitions. Brow-region signals were the strongest predictors under both SME and 3-Changepoint sampling, suggesting the classifiers converge on similar underlying signal and that the gap between strategies reflects frame quality rather than different learned features.

A further finding concerns reproducibility rather than accuracy. 3-Changepoint sampling reduced LSTM seed variance six-fold relative to the SME baseline, so sampling strategy affected the consistency of training even where it did not reliably affect final accuracy. Taken together, these results indicate that temporal sampling is a design decision with measurable consequences rather than a neutral preprocessing step, and one that argues against simply sampling more frames.

A reliable engagement classifier can inform agentic AI systems in educational settings, where detected engagement could trigger adaptive pacing, targeted hints, or personalised scaffolding. At roughly 70\% macro F1, however, such signals should serve as advisory inputs rather than deterministic triggers, so that interventions augment rather than override instructor judgement and learner agency.

Several limitations bound these conclusions and indicate the work needed next. The comparison covers two sampling families only, so uniform and random selection at a matched frame budget are needed before content-based selection can be preferred generally, and matched uniform segmentation would also permit direct comparison against the published EngageNet baselines \citep{singh2023}. Evaluation used a single binary threshold on a single dataset, leaving open whether the effect holds on the full four-level ordinal scale or transfers to corpora such as DAiSEE \citep{gupta2016}. Companion work examines how feature set and label definition interact on alternative datasets. The 3- to 7-frame budget also constrains sequence models by construction, so the absence of an LSTM advantage over the MLP should not be read as a general verdict on sequence modelling for this task, particularly since all architectures were compared at library defaults without hyperparameter search. Finally, the significance tests for the neural models rest on a single seed's predictions, and changepoint extraction via landmark velocity may not generalise to other salience signals used in motion-guided sampling.

% ---------------- REFERENCES ----------------
\begingroup
\fontsize{10}{12}\selectfont
\bibliography{references}
\endgroup

\end{document}
